% ============================================================================
% Hypothesis H5: climate-induced pressure on economic renewable potential
% relative to national expansion targets.
% H5a: resource-yield loss within the reference economic domain.
% H5b: exposure of the reference economic potential to climate hazards.
% Pre-registered decision rule (OQ-050, H5), fixed 2026-10-14, before any
% real-country result was produced. Materiality threshold lambda = 0.20.
% Notation follows docs/METHODOLOGY.md and the H1 to H4 sections.
% Citation keys marked "verify" must be checked before submission.
% ============================================================================

\section{Hypothesis H5: Climate-Induced Pressure on Economic Potential and
National Expansion Targets}
\label{sec:h5}

\subsection{Purpose and Position in the Research Design}
\label{subsec:h5_purpose}

Hypothesis H5 asks whether climate change puts material pressure on the
renewable potential that a country can use to meet its expansion target. H1 to
H4 concern which sites are best; H5 concerns how much economically viable
generation is at stake. It carries the link between climate risk and the energy
transition that motivates the thesis, and it is formulated so that a negative
answer is as informative as a positive one: if the pressure is small, climate
change does not threaten the feasibility of the target through the channels
represented.

Two mechanisms are tested separately.
\begin{itemize}
    \item H5a, resource. Climate change reduces the annual energy yield of the
    economically viable sites, through changes in irradiance, temperature and
    wind.
    \item H5b, exposure. A material part of the economically viable generation
    lies in cells where climate hazards exceed pre-specified physical
    thresholds.
\end{itemize}
A site may keep its yield and still be exposed to damaging conditions, and yield
may decline without hazard thresholds being exceeded; the two sub-hypotheses are
therefore evaluated independently before being combined.

Climate change affects energy systems both through gradual changes of the
resource and through changing extremes \cite{iaea2019}, and climate risk results
from the interaction of hazard, exposure and vulnerability
\cite{ipcc2022wgii}. H5b addresses exposure only. The temperature dependence of
photovoltaic performance \cite{skoplaki2009} and the projected changes of solar
and wind resources \cite{jerez2015,pryorbarthelmie2010} motivate H5a.
% verify jerez2015, pryorbarthelmie2010

\subsection{Economic Potential and Target Slack}
\label{subsec:h5_slack}

Renewable potential is defined at successive levels: technical potential is the
resource restricted by technology and land constraints, and economic potential
restricts it further by a cost criterion \cite{mercure2012,brown2015}. H5 uses
economic potential because a national target is met with competitive capacity,
not with every technically possible site. The levelised cost is a site-level
screening criterion; it does not represent system integration costs
\cite{ueckerdt2013}, and the economic potential of H5 is a screening volume, not
a deployment forecast.

\subsubsection{Reference economic domain}

For one country and one technology, let $\Omega$ be the set of candidate cells of
the central land scenario that have a valid climate forcing in every member of
the core window (the like-for-like set of the technical-potential phase). Cells
that become infeasible under some member are kept, since a loss of yield there is
part of what H5a measures. The reference economic domain is
\begin{equation}
    \Omega_{\mathrm{econ}}=\left\{\,i\in\Omega:\ \mathrm{LCOE}_{i,f_0}\leq\tau\,\right\},
    \label{eq:h5_domain}
\end{equation}
where $f_0=(m_0,s_0)$ is the nominal present future and $\tau$ is the levelised
cost threshold fixed in OQ-008, the same threshold that defines satisficing
robustness. The domain is defined once, at $f_0$, and held fixed across climate
members in the primary analysis.

\subsubsection{Energy and slack}

Let $E_{i,m}$ be the annual energy of cell $i$ under climate member $m$ in
TWh/year, computed by GeoFREA from the installable capacity of the central land
scenario and the capacity factor of the member; $E_{i,m_0}$ is the reference
value. The baseline economic potential is
\begin{equation}
    E_{0}=\sum_{i\in\Omega_{\mathrm{econ}}}E_{i,m_0}.
    \label{eq:h5_baseline}
\end{equation}
Let $G$ be the national target for the technology, in TWh/year, taken from the
official planning source fixed in OQ-010 and matched to the time horizon of the
core window. A target stated in installed capacity $G_{\mathrm{cap}}$ (GW) is
converted as
\begin{equation}
    G=8.76\,G_{\mathrm{cap}}\,\overline{\mathrm{CF}},
    \label{eq:h5_conversion}
\end{equation}
where $\overline{\mathrm{CF}}$ is the capacity-weighted reference capacity factor
of the cells of $\Omega_{\mathrm{econ}}$ taken in ascending order of
$\mathrm{LCOE}_{i,f_0}$ until their capacity reaches $G_{\mathrm{cap}}$. The
original unit and the conversion are reported with every result.

The target slack is the margin of the baseline economic potential over the
target:
\begin{equation}
    H_0=E_0-G .
    \label{eq:h5_slack}
\end{equation}
Both sub-hypotheses measure climate pressure relative to $H_0$. Normalising by
the total potential would make every effect negligible in countries whose
potential is many times their target; normalising by the slack asks how much of
the margin that supports the target is consumed or exposed. The ratios that
follow are measures of materiality, not probabilities of missing the target.

\subsection{H5a: Resource-Yield Loss}
\label{subsec:h5a}

For each member $m$ of the core window, the yield of the fixed domain is
$E_m=\sum_{i\in\Omega_{\mathrm{econ}}}E_{i,m}$. The loss statistic is the largest
reduction over members, truncated at zero:
\begin{equation}
    \Delta E=\max\!\left[\,0,\;E_0-\min_{m}E_m\,\right],
    \qquad
    R_{\mathrm{loss}}=\frac{\Delta E}{H_0}.
    \label{eq:h5a_statistic}
\end{equation}
The maximum over members follows the treatment of climate members throughout the
thesis: they are unweighted plausible futures, and a material adverse outcome in
any of them is relevant to planning. The truncation makes the test one-sided:
gains in some members do not offset losses in others. Gains, the ensemble median
and the member-specific values are reported as descriptive results.

Because the domain is fixed, $\Delta E$ measures the change in yield of the sites
that are economically viable today. It does not include sites that would enter or
leave the domain if the cost threshold were re-applied in each member; that
broader quantity is computed as a sensitivity.

\subsection{H5b: Exposure to Climate Hazards}
\label{subsec:h5b}

For each hazard channel $h$ with an indicator $Z_{i,h,m}$ and a threshold
$\theta_h$ fixed in OQ-051 (indicator, unit, temporal aggregation and comparison
rule), a cell is exposed under member $m$ when at least one channel reaches its
threshold:
\begin{equation}
    X_{i,m}=\mathbb{1}\!\left[\exists\,h:\ Z_{i,h,m}\geq\theta_h\right].
    \label{eq:h5_exposed_indicator}
\end{equation}
A cell exceeding several thresholds is counted once. The exposed economic energy
and the exposure ratio are
\begin{equation}
    E^{\mathrm{exp}}=\max_{m\in\mathcal{M}_{h}}
    \sum_{i\in\Omega_{\mathrm{econ}}}E_{i,m}\,X_{i,m},
    \qquad
    R_{\mathrm{exp}}=\frac{E^{\mathrm{exp}}}{H_0},
    \label{eq:h5b_statistic}
\end{equation}
where $\mathcal{M}_h$ is the set of members that carry the hazard channel: the
three climate models with a bias-adjusted hazard counterpart, in the core window.
The extreme-gust indicator is derived from historical reanalysis and does not
vary between members; it contributes identically to every member and to the
reference.

Exposure is not damage. A cell above a threshold is not assumed to fail or to
lose generation; that would require loss functions, which enter GeoFREA only when
supported by sufficient evidence. $R_{\mathrm{exp}}$ is a screening of how much of
the margin lies where hazards are severe. The same statistic evaluated with the
reference values of the indicators, and the difference between the two, are
reported to separate present exposure from the part added by climate change.

\subsection{Decision Rule and Verdict}
\label{subsec:h5_rule}

For each country and technology, with $\lambda=0.20$:
\begin{itemize}
    \item if no target $G$ is available, H5 is not testable; $E_0$, $\Delta E$ and
    $E^{\mathrm{exp}}$ are reported without ratios;
    \item if $H_0\leq0$, the status is target not met at reference; $\Delta E/G$
    and $E^{\mathrm{exp}}/G$ are reported as descriptive ratios, without a
    verdict;
    \item otherwise, H5a holds if $R_{\mathrm{loss}}\geq\lambda$ and H5b holds if
    $R_{\mathrm{exp}}\geq\lambda$, and the verdict is
\end{itemize}
\begin{equation}
    \text{H5 is}\;
    \begin{cases}
        \text{supported}, & \text{H5a and H5b hold},\\
        \text{partially supported (resource)}, & \text{only H5a holds},\\
        \text{partially supported (exposure)}, & \text{only H5b holds},\\
        \text{not supported}, & \text{neither holds}.
    \end{cases}
    \label{eq:h5_verdict}
\end{equation}
The two partial verdicts are kept distinct because they point to different
responses: a resource result to the margin over the target and the choice of
sites, an exposure result to design, adaptation and siting constraints.

The threshold $\lambda=0.20$ reads as one-fifth of the margin consumed or exposed,
the same proportion used for H1b and H3. It is a decision threshold of this
thesis; the literature motivates the comparison of climate pressure with planning
margins but provides no universal materiality cut-off. The defence rests on
pre-registration, on its reading and on the sensitivity reported below. The
inputs $\tau$, $G$ and $\theta_h$ come from sourced research and are fixed before
the results are produced; they are not adjusted to the outcome.

\subsection{Sensitivity Analyses}
\label{subsec:h5_sensitivity}

The verdict is issued from the primary configuration only; a change of verdict in
a variant is reported as a dependence on that choice.
\begin{enumerate}
    \item Materiality threshold $\lambda=0.10$, with the continuous ratios reported.
    \item Cost threshold at the ends of the sourced range of $\tau$.
    \item Dynamic economic domain: $\Omega_{\mathrm{econ},m}=\{i\in\Omega:
    \mathrm{LCOE}_{i,(m,s_0)}\leq\tau\}$, so that sites entering or leaving the
    domain under each member are included; reported separately from the
    fixed-domain loss.
    \item Land scenarios restrictive and permissive, which change the capacity of
    every cell and therefore $E_0$ and $H_0$.
    \item Climate: without IPSL-CM6A-LR; the 2071--2100 window for H5a (hazard
    indicators are not available for that window, and H5b is reported as not
    available); member-specific values and the identity of the members attaining
    the maxima.
    \item Hazard channels reported one by one, without summing overlapping cells.
    \item Target: alternative official targets or horizons where they exist, and
    alternative capacity-factor conversions for capacity targets.
    \item The 0.1-degree cell scale.
\end{enumerate}

\subsection{Worked Example}
\label{subsec:h5_example}

A country has $E_0=500$~TWh/year of economic potential and a target of
$G=300$~TWh/year, so $H_0=200$~TWh/year. The values are illustrative and are not
GeoFREA results.

H5a. The members give 480, 450 and 470~TWh/year over the fixed domain. Then
$\Delta E=50$~TWh/year and $R_{\mathrm{loss}}=0.25\geq0.20$: H5a holds. The most
adverse climate member consumes a quarter of the margin over the target.

H5b. The largest exposed economic energy among the hazard members is
60~TWh/year, so $R_{\mathrm{exp}}=0.30\geq0.20$: H5b holds. If the same statistic
with reference indicators is 40~TWh/year, climate change adds 20~TWh/year of
exposed potential.

Both hold, so H5 is supported for this country and technology. The verdict does
not mean that the target will be missed or that the exposed 60~TWh/year will be
lost: it means that, under the members and thresholds evaluated, both the yield
loss and the exposure reach one-fifth of the margin that supports the target.

\subsection{Parameters}
\label{subsec:h5_parameters}

\begin{table}[htbp]
    \centering
    \caption{Pre-registered parameters of H5.}
    \label{tab:h5_parameters}
    \small
    \renewcommand{\arraystretch}{1.2}
    \begin{tabularx}{\textwidth}{
        @{}
        >{\raggedright\arraybackslash}p{2.4cm}
        >{\raggedright\arraybackslash}p{3.2cm}
        >{\raggedright\arraybackslash}X
        @{}
    }
        \toprule
        \textbf{Parameter} & \textbf{Value} & \textbf{Role} \\
        \midrule
        $\lambda$ & 0.20 & Materiality threshold for H5a and H5b. \\
        \addlinespace
        $\tau$ & OQ-008 & Cost threshold of the economic domain; the same as
        satisficing robustness. \\
        \addlinespace
        $G$ & OQ-010 & National target in TWh/year, horizon-matched; capacity
        targets converted by Equation~\ref{eq:h5_conversion}. \\
        \addlinespace
        $\theta_h$ & OQ-051 & Hazard thresholds and indicator definitions. \\
        \addlinespace
        Domain & $\Omega_{\mathrm{econ}}$ at $f_0$ & Fixed across members in the
        primary analysis. \\
        \addlinespace
        Members & Core window; hazard members for H5b & Maximum over members,
        one-sided for losses. \\
        \bottomrule
    \end{tabularx}
\end{table}

\subsection{Limitations}
\label{subsec:h5_limitations}

Fixed domain. The primary loss concerns the sites that are economically viable
today; changes of eligibility under future members are covered only by the
dynamic-domain sensitivity.

Screening, not system adequacy. Annual energy does not represent seasonal
mismatch, multi-day low-production events, coincidence with peak demand or
network constraints. H5 is a national screening of yield and exposure, not a
reliability assessment.

Exposure, not risk. H5b measures exposure; vulnerability and consequences are not
modelled, and an exposed cell is not assumed to lose generation.

Hazard coverage. Exposure rests on three climate models in the core window and on
a historical gust indicator; hazards for the later window are not available.

Target definition. National targets usually refer to total installed capacity,
including existing plants, while the potential of GeoFREA does not remove the
cells already occupied by plants. The slack is therefore a screening margin, and
its value depends on the target source, horizon and unit conversion.

Climate representation. Climate enters through annual-mean change factors of the
resource, which may understate changes in yield; a not-supported H5a verdict is
conditional on this representation.

Generalisation. Results are conditional on the country, technology, cost
threshold, target, climate ensemble and hazard thresholds, and are reported per
country and technology.
